\section{Aggressive Cows Problem}
\label{sec:bs-aggressive-cows}

\problemstatement{We are given $n$ stall positions $\textit{stalls}[0..n-1]$
along a number line, and $k$ cows to place into these stalls (at most one
cow per stall). The cows are aggressive, so we want to place them to
\textbf{maximize the minimum distance} between any two cows. Find this
maximum possible minimum distance.}

\begin{patternbox}
\emph{``Maximize the minimum distance/gap between placed items''} is the
mirror image of every problem so far in this section: until now we
\emph{minimized} a candidate (speed, day, divisor, capacity) subject to a
feasibility constraint that became \emph{easier} as the candidate grew.
Here, feasibility (``can we place all $k$ cows with at least this much
spacing'') becomes \textbf{harder} as the candidate distance grows, so we
are finding the \textbf{largest} feasible candidate, not the smallest.
Whenever you see \emph{``maximize the minimum ...''}, expect to flip the
binary search's narrowing direction relative to the ``minimize subject to
feasibility'' problems seen earlier.
\end{patternbox}

\subsection*{Understanding the problem carefully}

Sort the stall positions first -- only relative order matters for spacing,
and an unsorted input gives no useful structure to exploit. For a fixed
candidate minimum distance $d$, ask: can we place all $k$ cows in the
sorted stalls such that every pair of adjacent placed cows is at least $d$
apart? This is checked greedily: place the first cow at the first stall,
then repeatedly place the next cow at the first stall that is at least $d$
away from the previously placed cow. If we manage to place all $k$ cows
this way, $d$ is feasible.

As $d$ increases, this greedy placement can only place fewer or equal
cows (every placement decision becomes at least as restrictive), so the
number of cows placeable is non-increasing in $d$ -- meaning feasibility
($\ge k$ cows placed) is \texttt{true} for small $d$ and \texttt{false}
for large $d$, the mirror image of the earlier sections' monotonicity
direction.

\begin{ideabox}[Candidate Space and Predicate]
The candidate space is distances $d \in [1, \max(\textit{stalls}) -
\min(\textit{stalls})]$. The predicate is
$P(d) = (\text{greedy placement with spacing} \ge d \text{ fits all } k
\text{ cows})$, and we want the \textbf{largest} $d$ where $P$ holds --
since $P$ is true for small $d$ and false for large $d$, we record a
successful candidate and narrow \emph{right} (search for an even larger
feasible $d$), the mirrored version of the ``narrow left after success''
template used for minimize-type problems.
\end{ideabox}

\subsection*{Why the greedy feasibility check is optimal}

For a \emph{fixed} target spacing $d$, greedily placing each next cow as
early as possible (at the first stall that is at least $d$ from the
previous placement) maximizes the number of cows we can fit using that
spacing: placing a cow any later than necessary only makes it harder to
fit subsequent cows, never easier, since stalls are processed in
increasing order and no stall can be revisited. This is the same kind of
exchange argument used for interval scheduling in the Greedy chapter --
the earliest valid placement is always at least as good as any later
valid placement for the purpose of fitting more items afterward.

\subsection*{Algorithm}

\begin{enumerate}
  \item Sort $\textit{stalls}$.
  \item Initialize $\textit{lo} \gets 1$,
        $\textit{hi} \gets \textit{stalls}[n-1] - \textit{stalls}[0]$.
  \item While $\textit{lo} \le \textit{hi}$:
  \begin{enumerate}
    \item $\textit{mid} \gets \textit{lo}+(\textit{hi}-\textit{lo})/2$.
    \item Greedily count how many cows fit with minimum spacing
          $\textit{mid}$.
    \item If $\textit{count} \ge k$: feasible; record \textit{mid}, try
          larger: $\textit{lo} \gets \textit{mid}+1$.
    \item Else: infeasible; try smaller: $\textit{hi} \gets \textit{mid}-1$.
  \end{enumerate}
  \item Return the recorded answer (equivalently \textit{hi} once the loop
        ends).
\end{enumerate}

\begin{algorithm}[H]
\caption{Aggressive Cows}
\begin{algorithmic}[1]
\Procedure{AggressiveCows}{\textit{stalls}, $k$}
  \State sort \textit{stalls}
  \State $\textit{lo} \gets 1$, $\textit{hi} \gets \textit{stalls}[n-1] - \textit{stalls}[0]$
  \While{$\textit{lo} \le \textit{hi}$}
    \State $\textit{mid} \gets \textit{lo}+(\textit{hi}-\textit{lo})/2$
    \State $\textit{count} \gets \Call{CowsPlaceable}{\textit{stalls}, \textit{mid}}$
    \If{$\textit{count} \ge k$}
      \State $\textit{lo} \gets \textit{mid} + 1$
    \Else
      \State $\textit{hi} \gets \textit{mid} - 1$
    \EndIf
  \EndWhile
  \State \Return \textit{hi}
\EndProcedure
\end{algorithmic}
\end{algorithm}

\subsection*{Dry run}

\dryrun{Take $\textit{stalls} = [0, 3, 4, 7, 10, 9]$, $k = 4$. Sorted:
$[0,3,4,7,9,10]$.}

\begin{itemize}
  \item $\textit{lo}=1,\textit{hi}=10$: $\textit{mid}=5$. Greedy: place at
        $0$; next $\ge 5$ is $7$, place; next $\ge 12$ -- none. Only $2$
        cows placed $< 4$. Infeasible. $\textit{hi}=4$.
  \item $\textit{lo}=1,\textit{hi}=4$: $\textit{mid}=2$. Greedy: place at
        $0$; next $\ge 2$ is $3$, place; next $\ge 5$ is $7$, place; next
        $\ge 9$ is $9$, place. $4$ cows placed $\ge 4$. Feasible.
        $\textit{lo}=3$.
  \item $\textit{lo}=3,\textit{hi}=4$: $\textit{mid}=3$. Greedy: place at
        $0$; next $\ge 3$ is $3$, place; next $\ge 6$ is $7$, place; next
        $\ge 10$ is $10$, place. $4$ cows placed $\ge 4$. Feasible.
        $\textit{lo}=4$.
  \item $\textit{lo}=4,\textit{hi}=4$: $\textit{mid}=4$. Greedy: place at
        $0$; next $\ge 4$ is $4$, place; next $\ge 8$ is $9$, place; next
        $\ge 13$ -- none. Only $3$ cows placed $< 4$. Infeasible.
        $\textit{hi}=3$.
  \item $\textit{lo}=4 > \textit{hi}=3$: loop ends.
\end{itemize}

The answer is $\textit{hi} = 3$: the maximum minimum distance is $3$.

\subsection*{C++ implementation}

\begin{lstlisting}
#include <bits/stdc++.h>
using namespace std;

int cowsPlaceable(vector<int>& stalls, int d) {
    int count = 1;
    int lastPos = stalls[0];

    for (int i = 1; i < (int)stalls.size(); i++) {
        if (stalls[i] - lastPos >= d) {
            count++;
            lastPos = stalls[i];
        }
    }
    return count;
}

int aggressiveCows(vector<int>& stalls, int k) {
    sort(stalls.begin(), stalls.end());
    int n = stalls.size();
    int lo = 1, hi = stalls[n - 1] - stalls[0];

    while (lo <= hi) {
        int mid = lo + (hi - lo) / 2;
        if (cowsPlaceable(stalls, mid) >= k) {
            lo = mid + 1;
        } else {
            hi = mid - 1;
        }
    }
    return hi;
}

int main() {
    vector<int> stalls = {0, 3, 4, 7, 10, 9};
    cout << "Maximum minimum distance: "
         << aggressiveCows(stalls, 4) << endl;
    return 0;
}
\end{lstlisting}

\begin{complexitybox}
\textbf{Time Complexity.} Sorting costs $O(n \log n)$. The candidate
distance range halves each iteration, costing
$O(\log(\max(\textit{stalls})-\min(\textit{stalls})))$ iterations, and
each feasibility check costs $O(n)$, giving an overall time complexity of
\[
  O(n \log n + n \log(\max(\textit{stalls}))).
\]
\textbf{Space Complexity.} $O(1)$ auxiliary space beyond sorting.
\end{complexitybox}

\begin{takeawaybox}
``Maximize the minimum'' and ``minimize the maximum'' problems both binary
search on the answer, but with the narrowing direction flipped relative to
each other: for ``maximize the minimum,'' record a successful candidate
and search for a \emph{larger} one; for ``minimize the maximum'' (as in
Sections~\ref{sec:bs-koko}--\ref{sec:bs-ship-packages}), record a
successful candidate and search for a \emph{smaller} one. Always pause to
identify which of the two flavors you are facing before writing the
narrowing logic.
\end{takeawaybox}
